% Track e5 -- The closure of the force law at two constituents (n = 2).
% Verification: investigations/emergence/verify/e5_closure_two.py (claims e5-1a ... e5-6b).
% New references: sections/bib_e5.tex.
\section{The closure at two constituents}
\label{e5:sec:main}

Section~17 of \cite{ES17} closes the force law $\dot p=\kappa G^{\!\top}p$ into the Lie--Poisson system of $\mathfrak s_n=\R\ltimes_{G_n}V_n$,
in which the rate is the angular velocity of a rotor. It treats one constituent, several constituents driven by one rotor
(Proposition~17.18 there), and three constituents (Theorems~17.19--17.24 there): the closure symmetry, the absence of a common Lorentz
algebra, the part of the Pati--Salam algebra that commutes with the binding root, what survives each projector, its chirality and
anomalies, and the freedom left in the inertia. The case $n=2$ is not treated there. Parts of it are in \cite{RE}: the closure
symmetry $\so{6,\C}$ (Proposition~7.8), the balance surface (Theorem~8.1), the balanced pairs (Proposition~8.4) and the pair invariant
$M_{\rm pair}$ (Theorem~8.6). This section supplies the $n=2$ analogue of every item, with the five Lorentz actions at $n=2$ in place of the two
projectors, which exist only at $n=3$, and asks whether $n=2$ is a closure rung in its own right. The answers, with their classes, are these. (1) The closure symmetry $\so{6,\C}$ is the complex orthogonal algebra of
$q_{\rm tot}=z_1\cdot z_1+z_2\cdot z_2$. It mixes the two field strengths, keeps no quadratic invariant except $q_{\rm tot}$, and does not
preserve the pair invariant: on each of its orbits $M_{\rm pair}$ takes every value in $[\half|q_{\rm tot}|,\infty)$ (fixed;
\cref{e5:sec:symmetry}). (2) There is no common Lorentz algebra, and under $\lG^{(1)}$ the spinor module carries no Weyl spinor
(fixed, partly restated; \cref{e5:sec:nocommon}). (3) There is no Pati--Salam algebra at $n=2$ (obstructed). The binding root leaves an
$\su2$ of the spectator gauge algebra, compact, non-abelian and kinetically admissible, which acts on the Weyl block with integer
isospin and vector-like; of the diagonal connection algebra it leaves the copy rotation; of the mixed rung it leaves nothing (fixed;
\cref{e5:sec:gauge}). (4) The inertia freedom is $(2,0)$ under the closure symmetry, as at $n=1$ and $n=3$, and it is $(14,10)$ or $(18,18)$
under the binding-compatible gauge parts (fixed, partly restated; \cref{e5:sec:inertia}). (5) The balance surface is connected and of
codimension two. It is transversally neutral, since the imbalance is a constant of motion; the rate is stable, but the bounded frame orbits
are not, and they need $h(0)\neq0$. With equal inertias a massive balanced pair is never an exchange eigenstate of its momenta
(fixed; \cref{e5:sec:rotor}).
(6) \Cref{e5:tab:ladder} compares $n=1,2,3$, and \cref{e5:sec:ladder} answers the question. Identifiers e5-$n$ refer to the \texttt{PASS}
lines of \texttt{verify/e5\_closure\_two.py} (\cref{e5:tab:claims}).

\paragraph{Notation.} As in Section~8 of \cite{RE}. $V_2=V\oplus V$ has the basis $e^{(1)}_a$, $e^{(2)}_a$, the metric $\eta_2=\eta\oplus\eta$ and
the binding root $G_2=G\oplus G$. The temporal directions are $e^{(k)}_1,e^{(k)}_2,e^{(k)}_3$, each with its partner $e^{(k)}_{i+3}$, and the
paired rotations and boosts $\mathrm{rot}(i,j)$, $\mathrm{bst}(i,j)$ of temporal $i<j$ are those of Section~17.6 of \cite{ES17}. A vector $X$ of
$V_2$ gives $Z=(z_1,z_2)\in\C^6$, $z_k=z(X^{(k)})$, on which $G_2$ acts as $i$; $q_k=z_k\cdot z_k$, $q_{\rm tot}=q_1+q_2$, $Q_{kl}=z_k\cdot z_l$ is the
complex Gram matrix with singular values $\sigma_1\ge\sigma_2$, and $M_{\rm pair}=\half(\sigma_1+\sigma_2)$ is the rest energy of the pair
(Theorem~8.6 of \cite{RE}). $C_v=\varepsilon\otimes\id_6$ is the copy rotation and $\varepsilon\otimes G$ the copy boost,
$\varepsilon=\bigl(\begin{smallmatrix}0&-1\\1&0\end{smallmatrix}\bigr)$. $\lG^{(k)}$ and $\lint^{(k)}$ are the frame and interior Lorentz
algebras of constituent $k$, and $\lG^\Delta$ is the diagonal $\lG$. On $S_2=\R^{64}$ the tower generators are $E_{(1,a)}=e_a\otimes\id$,
$E_{(2,a)}=\omega\otimes e_a$, the spin lift of $M\in\so{6,6}$ is $S(M)=\tfrac14\sum_{a,b}(M\eta_2)_{ab}E_aE_b$, and $\Gamma_{12}=E_{(1,1)}\cdots E_{(2,6)}$.
In the closure the per-constituent quantities are those of the momentum fields $\pi^{(k)}=\eta^{-1}p^{(k)}$, $w_k=w(p^{(k)})=q(z(\pi^{(k)}))$
(equation~(30) of \cite{ES17}); Lorentz-class inertias are $N_k=(a_k\eta^{-1}+b_kG\eta^{-1})/m_k$, $c_k=(a_k+ib_k)/m_k$, $F(p)=\sum_kc_kw_k$ and
$R_{\rm eff}e^{i\varphi_{\rm eff}}=F(p(0))$. The \emph{kinetic condition} is that of Proposition~5.5 of \cite{ES17}, used in Theorem~17.22 there
with the kernels compressed by the projector: $X^{\!\top}K^\mu+K^\mu X=0$ for $\mu=0,\dots,3$, with the kernels $K^\mu=-\gamma^0\gamma^\mu$ of
Proposition~4.5 there. At $\mu=0$ it says that $X$ is antisymmetric, and at $\mu=a$ that $X$ commutes with the boosts $\gamma^0\gamma^a$.

\subsection{The closure symmetry and the invariants of a pair}
\label{e5:sec:symmetry}

\begin{proposition}[The closure symmetry at two constituents]
\label{e5:prop:so6C}
$G_2^{\!\top}\eta_2G_2=-\eta_2$ and $G_2^2=-\id$. The centraliser of $G_2$ in $\so{6,6}$ has dimension $30$ and is spanned by the $15$ paired
rotations and the $15$ paired boosts. In the basis of the six temporal directions followed by their partners every element is
$X=\bigl(\begin{smallmatrix}R&S\\-S&R\end{smallmatrix}\bigr)$ with $R,S$ antisymmetric, and $X$ acts on $Z$ as the complex antisymmetric matrix
$R-iS$. The map $X\mapsto R-iS$ is a Lie isomorphism onto $\so{6,\C}$, the algebra of the complex orthogonal group of $q_{\rm tot}$, and the
compact part of the centraliser is $\so6$, the paired rotations.
\end{proposition}

This restates Proposition~7.8(i) of \cite{RE} (fixed there), the $n=2$ member of Theorem~17.19 of \cite{ES17}; the bracket identity was checked
on all $435$ pairs of basis elements (e5-1a). By Proposition~7.8(ii) of \cite{RE}, $\so{6,\C}=\lG^{(1)}\oplus\lG^{(2)}\oplus(\C^3\otimes\C^3)$,
and the cross part maps $z_1$ into $z_2$: the closure symmetry does not keep the two field strengths apart.

\begin{proposition}[Quadratic invariants of a pair]
\label{e5:prop:quad}
The holomorphic quadratic polynomials in $(z_1,z_2)$ invariant under
(a) $\so{6,\C}$ are the multiples of $q_{\rm tot}$;
(b) $\lG^{(1)}\oplus\lG^{(2)}$ form $\operatorname{span}_\C\{q_1,q_2\}$;
(c) $\lG^\Delta$ form $\operatorname{span}_\C\{q_1,q_2,z_1\cdot z_2\}$.
In each case the invariant real quadratic forms on $V_2$ have real dimension $2$, $4$, $6$ and are spanned by the real and imaginary parts of these.
\end{proposition}

\begin{proof}
$\lG\cong\so{3,\C}$ acts on $\C^3$ by the vector representation. By the first fundamental theorem for the orthogonal group \cite{Weyl39},
the polynomial invariants of $SO(3,\C)$ acting diagonally on $(z_1,z_2)$ are generated by the inner products $z_k\cdot z_l$ (and
$\det(z_1,z_2,\cdot)$, which needs three vectors), which gives (c). Under $\lG^{(1)}\oplus\lG^{(2)}$ a cross term $z_1\cdot Az_2$ is invariant
under $z_1\mapsto\Lambda z_1$ with $z_2$ fixed only if $A=0$, which gives (b). $\C^6$ is irreducible under $\so{6,\C}$, with $q_{\rm tot}$ its
only invariant form, which gives (a). A real invariant form that is not the real part of a holomorphic one would be Hermitian, $h(Xu,v)+h(u,Xv)=0$,
and applied to $X$ and $iX$ this forces $h(Xu,v)=0$ for all $X$, so $h=0$.
\end{proof}

Inherited (classical invariant theory), with the real dimensions computed as exact null spaces (e5-1b).

\begin{theorem}[The closure symmetry and the pair invariant]
\label{e5:thm:pair}
\begin{enumerate}[label=(\roman*),leftmargin=2.2em]
\item $q(z_1+z_2)=q_1+q_2+2\,z_1\cdot z_2$, $\det Q=q_1q_2-(z_1\cdot z_2)^2$, and
\[
M_{\rm pair}^2=\tfrac14\bigl(|q_1|^2+|q_2|^2+2|z_1\cdot z_2|^2\bigr)+\tfrac12\bigl|q_1q_2-(z_1\cdot z_2)^2\bigr| .
\]
So $M_{\rm pair}$, $z_1\cdot z_2$ and $q(z_1+z_2)$ are Lorentz invariants of the pair under $\lG^\Delta$. None of them is invariant under
$\lG^{(1)}\oplus\lG^{(2)}$: a boost of constituent~$2$ alone fixes $q_1,q_2$ and changes $z_1\cdot z_2$ and $M_{\rm pair}$. Under the closure
symmetry only $q_{\rm tot}=\tr Q$ is invariant.
\item $M_{\rm pair}$ is invariant under $\lG^\Delta$ and the copy rotation $C_v$, and the subalgebra of $\so{6,\C}$ that leaves it invariant
is $\lG^\Delta\oplus\R C_v$, of real dimension $7$.
\item The closure symmetry does not preserve $M_{\rm pair}$. The copy boost has the eigenvector $(e_1,ie_1)$ with eigenvalue $1$, and maps the
balanced pair $(e_1,ie_1)$, with $q_{\rm tot}=0$ and $M_{\rm pair}=1$, to $e^t(e_1,ie_1)$, with $M_{\rm pair}=e^{2t}$. The compact rotation
$R^{(1)}_\varphi\in SO(6)$ of Example~8.8 of \cite{RE} takes $((1,i,0),0)$ from $M_{\rm pair}=0$ to $M_{\rm pair}=\sin\varphi$.
\item For each $c\in\C$ the set $\{Z\neq0:\ q_{\rm tot}(Z)=c\}$ is a single orbit of $SO(6,\C)$, and on it $M_{\rm pair}$ takes every value in
$[\half|c|,\infty)$. The bound $M_{\rm pair}\ge\half|q_{\rm tot}|$ of Theorem~8.6(iii) of \cite{RE} is attained at $(z,0)$, and the copy boost
carries $(z,0)$ to $(\cosh t\,z,\,i\sinh t\,z)$, with $q_{\rm tot}=q(z)$ and $M_{\rm pair}=\half|q(z)|\cosh2t$.
\end{enumerate}
\end{theorem}

\begin{proof}
(i) The first two identities are expansions, and with them the formula is $M_{\rm pair}^2=\tfrac14\|Q\|_F^2+\half|\det Q|$, which follows from
Theorem~8.6(ii) of \cite{RE}. A diagonal Lorentz transformation $\Lambda$ maps $Q=Z^{\!\top}Z$ to $Z^{\!\top}\Lambda^{\!\top}\Lambda Z=Q$; one acting on
constituent~$2$ alone maps $z_1\cdot z_2$ to $z_1\cdot\Lambda z_2$. The last statement is \cref{e5:prop:quad}(a). (ii) $C_v$ maps $Q$ to $RQR^{\!\top}$ with
$R$ real orthogonal, which keeps the singular values (Proposition~8.9(ii) of \cite{RE}). That nothing else in $\so{6,\C}$ keeps $M_{\rm pair}$ is
computed: the derivatives of $M_{\rm pair}^2$ along the $30$ basis fields at $40$ random pairs form a matrix of numerical rank $23$. (iii) The copy
boost acts on the copy index as $\bigl(\begin{smallmatrix}0&-i\\i&0\end{smallmatrix}\bigr)$ (Proposition~7.8(iii) of \cite{RE}), and $M_{\rm pair}$
is homogeneous of degree two. At $(e_1,ie_1)$, $Q=vv^{\!\top}$ with $v=(1,i)$, so $\sigma_1=2$ and $\sigma_2=0$. The numbers for $R^{(1)}_\varphi$
restate Example~8.8 of \cite{RE}. (iv) Transitivity on the level sets of a nondegenerate quadratic form is Witt's theorem \cite{Witt37}. At $(z,0)$,
$Q=\operatorname{diag}(q(z),0)$. Since $(\varepsilon\otimes G)^2=\id$, $e^{t\,\varepsilon\otimes G}=\cosh t+\sinh t\,\varepsilon\otimes G$, and the image has
$Q=q(z)\,vv^{\!\top}$ with $v=(\cosh t,i\sinh t)$, so $\sigma_1=|q(z)|\,|v|^2=|q(z)|\cosh2t$ and $\sigma_2=0$. The orbit is connected and $M_{\rm pair}$
is continuous, so it takes every value between its infimum and its unbounded values; for $c=0$ use (iii).
\end{proof}

Items (i) and (iv) are fixed, with (i) partly restated from Theorem~8.6 of \cite{RE}; (ii) is fixed, with the stabiliser computed numerically;
(iii) is fixed, with the second example restated. Checks: the identities symbolically, the formula against the singular values at $30$ random
pairs to $10^{-12}$ (e5-1e); invariance to $10^{-11}$, and the rank $23$ with a singular-value gap above $10^8$ (e5-1c); the orbits to
$10^{-12}$ (e5-1d). So the closure symmetry mixes the two field strengths and does not preserve the pair invariant. The closure symmetry sees one field strength in $\C^6$ and its square $q_{\rm tot}$. To split it into two
field strengths that carry a pair invariant takes the seven-dimensional stabiliser $\lG^\Delta\oplus\R C_v$, a diagonal Lorentz action together
with the copy rotation, and the closure selects neither. With equal inertias its dynamics agrees: the frame energy is $\half\Re(c\,q_{\rm tot}(\pi))$
(Proposition~17.18 of \cite{ES17}; \cref{e5:prop:rotorlaw}(i) with $c_1=c_2=c$), and its compact symmetry carries every balanced pair to one null constituent and an empty copy (Proposition~8.9(iii)
of \cite{RE}). In that closure ``two constituents'' is a label that the dynamics does not see (interpretive).

\subsection{No common Lorentz algebra}
\label{e5:sec:nocommon}

\begin{proposition}[No common Lorentz algebra at two constituents]
\label{e5:prop:nocommon}
\begin{enumerate}[label=(\roman*),leftmargin=2.2em]
\item $\Cent(G_2)\cap\lint^{(1)}=0$. The elements of $\so{6,6}$ commuting with both $G_2$ and $\lint^{(1)}$ form $\lG^{(2)}\cong\so{3,\C}$, of
dimension $6$. At rung $n$ they form the $\so{3n-3,\C}$ of constituents $2,\dots,n$, of dimension $3(n-1)(3n-4)$: $0$, $6$ and $30$ at
$n=1,2,3$.
\item The bivectors of constituent~$1$ are $e_ae_b\otimes\id$, so $S_n\otimes\C$ carries $8^{n-1}$ copies of the content of $S\otimes\C$. Under
$\lG^{(1)}$, $S_2\otimes\C=16\,(\half,\half)$: no Weyl spinor ($128\,(\half,\half)$ at $n=3$). Under $\lint^{(1)}$,
$S_2\otimes\C=16\,(\half,0)\oplus16\,(0,\half)$.
\item Under $\lG^{(2)}$, $S_2\otimes\C=16\,(\half,\half)$, and under $\lG^\Delta$,
$S_2\otimes\C=4\bigl[(0,0)\oplus(1,0)\oplus(0,1)\oplus(1,1)\bigr]$: integer spin only.
\end{enumerate}
\end{proposition}

\begin{proof}
(i) An element of $\lint^{(1)}$ lives on constituent~$1$, where $\Cent(G)\cap\so{3,3}=\lG\cong\so{3,\C}$ (Theorem~12.1 of \cite{ES17}). It fixes the
surplus pair, and an element of $\so{3,\C}$ that kills $z=(0,1,0)$ and $z=(0,0,1)$ is an antisymmetric matrix with two zero columns, hence zero. Under
$\lint^{(1)}$, $V_2$ is the interior four-plane of constituent~$1$, irreducible and nontrivial, plus eight Lorentz-trivial directions: the surplus
pair $F=\langle e^{(1)}_2,e^{(1)}_3\rangle$ and all of constituent~$2$. By Schur's lemma an equivariant $X\in\so{6,6}$ vanishes on the four-plane
and preserves the eight trivial directions. If $X$ also commutes with $G_2$, then for $v\in F$, $G_2v$ lies in the interior plane
$\langle e^{(1)}_5,e^{(1)}_6\rangle$, so $G_2Xv=XG_2v=0$ and $Xv=0$. Hence $X$ vanishes on constituent~$1$ and is an element of $\so{3,3}$ of
constituent~$2$ commuting with $G$, that is of $\lG^{(2)}$ (Theorem~12.1 of \cite{ES17}). The same argument at rung $n$ leaves
$\Cent(G_{n-1})$ of constituents $2,\dots,n$, which is $\so{3n-3,\C}$ by Proposition~7.8(i) of \cite{RE}. (ii) $E_{(1,a)}E_{(1,b)}=e_ae_b\otimes\id$,
and at one constituent $S\otimes\C=2\,(\half,\half)$ under $\lG$ and $2\,(\half,0)\oplus2\,(0,\half)$ under $\lint$ (Section~13 of \cite{ES17}).
(iii) The bivectors of constituent~$2$ are $\id\otimes e_ae_b$, so $\lG^\Delta$ acts on $(S\otimes\C)\otimes_\C(S\otimes\C)$, and
$(\half,\half)\otimes(\half,\half)=(0,0)\oplus(1,0)\oplus(0,1)\oplus(1,1)$.
\end{proof}

Part (i) restates Proposition~7.8(iv) and Proposition~9.1(iii) of \cite{RE} (fixed there), the dimension $30$ restates Theorem~17.20 of
\cite{ES17}, and the ladder formula is fixed (e5-2a). Parts (ii) and (iii) are fixed at $n=2$, restated at $n=3$, and the Clebsch--Gordan rule is
inherited (Casimir spectra to $10^{-9}$, computed on $S_2$ and on one constituent; e5-2b, e5-2c). No nonzero Lorentz generator leaves both the
closure and the spinor structure of matter invariant, at $n=2$ as at $n=1$ and $n=3$. What differs is what commutes with both: at $n=2$ it is a
whole second Lorentz algebra, the frame algebra of the spectator constituent. That is the material of the mixed rung, which pairs $\lint^{(1)}$ with
$\lG^{(2)}$ (\cref{e5:prop:mixed}).

\subsection{Gauge algebras at two constituents and the binding root}
\label{e5:sec:gauge}

\begin{proposition}[No Pati--Salam algebra at two constituents]
\label{e5:prop:noPS}
The compact parts of the candidate gauge algebras $\mathfrak g(k_I,k_B,k_S)$ of the five Lorentz actions at $n=2$ (Theorem~10.8 of \cite{ES17};
Table~8 of \cite{RE}) have dimensions $13$ (spectator, $\so5\oplus\so3$), $6$ (binding and spectator, $\so3\oplus\so3$), $7$ (diagonal interior,
$\so2\oplus\so4$), $1$ (mixed) and $1$ (binding on both). None contains the Pati--Salam algebra $\so4\oplus\so6$ \cite{PatiSalam}, of dimension $21$,
or its colour factor $\su4$, of dimension $15$.
\end{proposition}

Obstructed, by a dimension count; the five dimensions are re-verified (e5-3a, e5-3e, e5-3h, e5-3i). At $n=3$ the compact spectator algebra
$\so8\oplus\so6$, of dimension $43$, already contains $\mathfrak{ps}_U$ (Section~10.7 of \cite{ES17}). The analogue at $n=2$ of
Theorems~17.21--17.23 of \cite{ES17} is therefore the intersection of each compact candidate with $\Cent(G_2)$.

\begin{proposition}[The spectator algebra and the binding root]
\label{e5:prop:spectator}
Let the Lorentz algebra act by $\lint^{(1)}$, with constituent~$2$ a spectator.
\begin{enumerate}[label=(\roman*),leftmargin=2.2em]
\item $\mathfrak g(1,0,1)=\Cent(\lint^{(1)})\cap\so{6,6}\cong\so{5,3}$ has dimension $28$. Its compact part is $\so5\oplus\so3$, with $\so5$ on the
positive Lorentz-trivial span $P=\langle e^{(1)}_2,e^{(1)}_3,e^{(2)}_1,e^{(2)}_2,e^{(2)}_3\rangle$ and $\so3$ on $N=\langle e^{(2)}_4,e^{(2)}_5,e^{(2)}_6\rangle$.
$\mathfrak g(1,0,1)\cap\Cent(G_2)=\lG^{(2)}$. Neither $\so5$ nor $\so3$ meets $\Cent(G_2)$, and with $E_2=\langle e^{(2)}_1,e^{(2)}_2,e^{(2)}_3\rangle$,
\[
(\so5\oplus\so3)\cap\Cent(G_2)=\bigl\{A+G_2AG_2^{-1}:\ A\in\so{E_2}\bigr\}\cong\su2 ,
\]
the three paired rotations of constituent~$2$: compact and non-abelian.
\item Among the spin lifts of the $30$ elements of $\Cent(G_2)$, those satisfying the kinetic condition with the kernels of constituent~$1$ are
exactly the paired rotations of constituent~$2$. The paired boosts of constituent~$2$ fail it.
\item Let $w$ be the product of the eight Lorentz-trivial generators. Then $w^2=-\id$, $w$ commutes with $\lint^{(1)}$ and with the spin lift of
$\so{5,3}$, and $\Gamma=w\,\Omega^{(1)}$, with $\Omega^{(1)}=E_{(1,1)}E_{(1,4)}E_{(1,5)}E_{(1,6)}$ the interior volume element of
constituent~1, is a chirality commuting with both. With $w$ as $i$, each chirality half is a Weyl block $(\half,0)\otimes(\mathbf4,\mathbf2)$, respectively $(0,\half)\otimes(\mathbf4,\mathbf2)$, of
$\lint^{(1)}\oplus\so5\oplus\so3$ (Proposition~10.12 of \cite{ES17}). Under the binding-compatible $\su2$,
\[
L=(\half,0)\otimes\bigl(2\cdot\mathbf1\oplus2\cdot\mathbf3\bigr),\qquad R=(0,\half)\otimes\bigl(2\cdot\mathbf1\oplus2\cdot\mathbf3\bigr):
\]
integer isospin, and the same real representation on both halves, so the action is vector-like. With charges $h=-w\,S(X)$,
$\mathcal A(h)=\mathcal A(h^3)=0$, where $\mathcal A(f)=\tfrac14[\Tr(fP_L)-\Tr(fP_R)]$ as in Theorem~17.23 of \cite{ES17}. There is no doublet, hence no
global $SU(2)$ anomaly \cite{Witten82}.
\end{enumerate}
\end{proposition}

\begin{proof}
(i) Let $A\in\so P$ and $B\in\so N$ with $[A+B,G_2]=0$, and follow the proof of Theorem~17.21 of \cite{ES17}. $G_2$ maps $F=\langle e^{(1)}_2,e^{(1)}_3\rangle$
into the interior plane $\langle e^{(1)}_5,e^{(1)}_6\rangle$, on which $A$ and $B$ vanish, so $G_2Av=(A+B)G_2v=0$ and $Av=0$ for $v\in F$. For $v\in E_2$,
$Bv=0$ and $G_2v\in N$, so $G_2Av=BG_2v\in N$ and $Av\in G_2^{-1}N=E_2$. Hence $A\in\so{E_2}$ and $B=G_2AG_2^{-1}$ on $N$; conversely every such pair
commutes with $G_2$. Taking $B=0$ or $A=0$ shows that neither summand meets $\Cent(G_2)$. The rest is \cref{e5:prop:nocommon}(i). (ii) At $\mu=0$ the
condition is antisymmetry. The spin lift of a paired boost is a sum of boost bivectors, which are symmetric (Theorem~2.3 of \cite{ES17}). The spin lift
of a paired rotation of constituent~$2$ is antisymmetric and commutes with every bivector of constituent~$1$, in particular with $\gamma^0\gamma^a$.
A general element of $\Cent(G_2)$ was tested by an exact null space. (iii) $\so5\supset\so{E_2}\oplus\so F$ restricts the spinor $\mathbf4$ to two
isospin doublets, and $\so3$ acts on $\mathbf2$ as a doublet, so the diagonal $\su2$ sees $(\mathbf2\oplus\mathbf2)\otimes\mathbf2=2\cdot(\mathbf1\oplus\mathbf3)$.
Every representation of $\su2$ has weights symmetric under $m\mapsto-m$, so every odd trace of $h$ vanishes on each half
separately; integer isospin excludes doublets, and with them the global anomaly.
\end{proof}

Part (i) is fixed: it is the structure of Theorem~17.21 of \cite{ES17} at $n=2$, with $U=P$ and $d=\dim E_2=3$. Part (ii) is fixed, and in
(iii) the Weyl blocks restate Proposition~10.12 of \cite{ES17}, the branching is inherited, and the content and the anomaly traces are fixed. Checks:
exact Lie algebra (e5-3a), exact null space of the kinetic condition (e5-3d), Casimirs to $10^{-9}$ and integer traces (e5-3b, e5-3c); the odd traces vanish as well for the $\so3$ on $N$, which
acts by doublets (e5-3j). The binding
root therefore turns the isospin doublets of the spectator algebra, the ``doublets, no colour, real'' of Section~10.8 of \cite{ES17}, into isospin
$0$ and $1$. No doublet survives at $n=2$, as for the copy isospin of Proposition~9.6 of \cite{RE}. At $n=3$ the doublets survive in the
$(\half,\half)_{\pm1}$ of $\so4_\Delta$ on $P_*$ (Theorem~17.23(ii) of \cite{ES17}).

\begin{proposition}[The diagonal connection algebra and the binding root]
\label{e5:prop:diagonal}
Let the Lorentz algebra act by the diagonal $\lint$.
\begin{enumerate}[label=(\roman*),leftmargin=2.2em]
\item $\Cent(\text{diagonal }\lint)\cap\so{6,6}=\so2\oplus\so4$ has dimension $7$ and is compact. It meets $\Cent(G_2)$ in $\R C_v$, the copy
rotation of all six directions: the copy rotation of the interior four-planes plus that of the surplus pairs. This algebra is abelian and compact,
and $S(C_v)=-C_b$ with $C_b=\half\sum_a\eta_{aa}E_{(1,a)}E_{(2,a)}$.
\item There is no fermion under this action at $n=2$ (Proposition~10.9 of \cite{ES17}) and no copy-invariant gamma matrix (Theorem~13.15 of
\cite{ES17}), so the kinetic condition has no canonical kernel. With the summed kernels $K^\mu_S=\sum_k\gamma^{(k)}_0\gamma^{(k)}_\mu$ it holds for
$C_b$. With the kernels $f_0f_\mu$ of the centre-of-mass generators $f_a=(E_{(1,a)}+E_{(2,a)})/\sqrt2$ (Proposition~7.3 of \cite{RE}) it fails.
\item $\Gamma_{12}^2=\id$, and $\Gamma_{12}$ commutes with the diagonal $\lint$ and with $C_b$. With copy charge $q$ defined by $C_b=iq$, the content
of $S_2\otimes\C$ is
\begin{center}
\small
\begin{tabular}{@{}ll@{}}
\toprule
$\Gamma_{12}=-1$ & $(0,0)_{\pm3}\oplus(0,0)_{\pm1}\oplus2\,(\half,\half)_{\pm1}\oplus(1,0)_{\pm1}\oplus(0,1)_{\pm1}$\\
$\Gamma_{12}=+1$ & $2\,(0,0)_{\pm2}\oplus(\half,\half)_{\pm2}\oplus2\,(\half,\half)_0\oplus2\,(1,0)_0\oplus2\,(0,1)_0$\\
\bottomrule
\end{tabular}
\end{center}
Spin is integral, $\Gamma_{12}=-1$ carries the odd and $\Gamma_{12}=+1$ the even copy charges, and each half is invariant under
$(q,(j_L,j_R))\mapsto(-q,(j_R,j_L))$: self-conjugate, hence vector-like. $\Tr(C_b^kP_\pm)=0$ for every odd $k$.
\end{enumerate}
\end{proposition}

\begin{proof}
(i) Under the diagonal $\lint$, $V_2=2V\oplus4\cdot\mathbf1$, and Schur's lemma gives $\so2\oplus\so4$, the $n=2$ case of Theorem~10.2 of
\cite{ES17}. Within it, an element $aM+B$, with $M$ the copy rotation of the interior four-planes and $B$ in the surplus $\so4$, commutes with
$G_2$ exactly when $B$ is $a$ times the copy rotation of the surplus pairs, because $G_2$ maps each surplus pair into the interior plane of the same
constituent. This is Proposition~9.1(ii) of \cite{RE}, and the lift is Proposition~7.3(ii) there. (ii) $C_b$ commutes with the diagonal boosts and is
antisymmetric, while $f_0f_\mu$ contains the cross terms $E_{(1,1)}E_{(2,\mu)}$ ($\gamma^0=e_1$), with which it does not commute. (iii) Computed.
\end{proof}

Part (i) restates Proposition~9.1(ii) and Proposition~7.3 of \cite{RE} (fixed there; e5-3e), (ii) is fixed (exact; e5-3f), and (iii) is fixed
(eigenvalues to $10^{-9}$, traces exact; e5-3g). Under the diagonal action the binding root leaves an abelian, compact, Lorentz-invariant copy charge
on matter that has no fermions; the row ``$\so2\oplus\so4$, no fermions'' of Section~10.8 of \cite{ES17} becomes ``$\so2$, no fermions''.

\begin{proposition}[The mixed rung and the binding root]
\label{e5:prop:mixed}
Let $\mathfrak l_{\rm mix}$ pair $\lint^{(1)}$ with $\lG^{(2)}$ generator by generator (Proposition~13.16 of \cite{ES17}).
\begin{enumerate}[label=(\roman*),leftmargin=2.2em]
\item $\mathfrak l_{\rm mix}$ satisfies the $\so{1,3}$ brackets, its centraliser in $\so{6,6}$ is $\R J_1$, $J_1=e^{(1)}_2e^{(1)}_3$ acting on vectors
(Proposition~13.17 of \cite{ES17}), and $\R J_1\cap\Cent(G_2)=0$. $G_2$ does not commute with $\mathfrak l_{\rm mix}$.
\item The $\mathfrak l_{\rm mix}$-equivariant endomorphisms of $V_2$ form a seven-dimensional algebra: the scalars on the four-vector, $\operatorname{span}\{\id,G\}$ on
the field strength of constituent~$2$, and $\mathfrak{gl}(2)$ on the two Lorentz-trivial directions. Its Euclidean-antisymmetric part, the generators
that rotate without dilating, is $\operatorname{span}\{G^{(2)},J_1\}$, with $G^{(2)}=0\oplus G$, and $[J_1,G^{(2)}]=0$.
\end{enumerate}
\end{proposition}

\begin{proof}
(i) $\R^{6,6}=V\oplus\Lambda^2V\oplus2\cdot\mathbf1$ under $\mathfrak l_{\rm mix}$ (Proposition~13.16 of \cite{ES17}). $J_1$ rotates the surplus pair of
constituent~$1$, and $G_2$ maps that pair into the interior four-plane, which $J_1$ fixes, so $[J_1,G_2]\neq0$. No nonzero element of
$\lint^{(1)}$ commutes with $G_2$ (\cref{e5:prop:nocommon}(i)), so $G_2$ does not commute with $\mathfrak l_{\rm mix}$. (ii) Schur's lemma on the three isotypic pieces, with $\Hom_{\lG}(V,V)=\operatorname{span}\{\id,G\}$
(Proposition~17.4 of \cite{ES17}).
\end{proof}

Part (i) restates Propositions~13.16 and~13.17 of \cite{ES17} and is fixed in what it adds; (ii) is fixed (exact; e5-3h). The binding root of the
closure is not covariant under the mixed action and leaves nothing of its algebra. The only covariant rotor generators that rotate without
dilating are the combinations of $G^{(2)}$ and $J_1$. With the rotor $G^{(2)}$, which binds constituent~$2$ alone, all of $\R J_1$ survives
(constructed: admitted, not selected). For completeness, under the two binding actions the binding root leaves $\lG^{(2)}$ of
$\mathfrak g(0,1,1)=\so{3,3}^{(2)}$, with compact part the $\su2$ of \cref{e5:prop:spectator}, and all of $\mathfrak g(0,2,0)=\so{2,\C}$, with compact part
$\R C_v$ (restating Proposition~7.8(iii) and Table~8 of \cite{RE}; e5-3i). Neither action carries a Dirac operator, so the kinetic condition does
not apply.

\subsection{Inertia freedom at two constituents}
\label{e5:sec:inertia}

\begin{theorem}[Inertia freedom at two constituents]
\label{e5:thm:inertia}
The real dimensions of the symmetric forms $N$ on $V_2^*$, invariant under a symmetry ($rN+Nr^{\!\top}=0$, Definition~17.10 of \cite{ES17}), in the
Lorentz class $\{N,G_2\}=0$ and in the duality class $[N,G_2]=0$, are:
\begin{center}
\small
\begin{tabular}{@{}lcc@{}}
\toprule
symmetry & Lorentz class & duality class\\
\midrule
$\Cent(G_2)=\so{6,\C}$ & $2$ & $0$\\
$\lG^\Delta\oplus\R C_v$ (the stabiliser of $M_{\rm pair}$) & $2$ & $0$\\
$\lG^{(1)}\oplus\lG^{(2)}$ & $4$ & $0$\\
$\lG^\Delta$ & $6$ & $0$\\
$\lG^{(2)}$ & $14$ & $9$\\
$\su2$, the binding-compatible spectator part (\cref{e5:prop:spectator}) & $14$ & $10$\\
$\R C_v$, the binding-compatible diagonal part (\cref{e5:prop:diagonal}) & $18$ & $18$\\
none & $42$ & $36$\\
\bottomrule
\end{tabular}
\end{center}
Under the full closure symmetry $\so{3n,\C}$ the freedom is $(2,0)$ at $n=1,2,3$, spanned by $\eta_n^{-1}$ and $G_n\eta_n^{-1}$; with no symmetry it is
$(3n(3n+1),9n^2)$, that is $(12,9)$, $(42,36)$ and $(90,81)$.
\end{theorem}

\begin{proof}
Split $N$ into blocks between constituents. Under $\lG^{(k)}$ the invariant symmetric forms of one constituent are $\operatorname{span}\{\eta^{-1},G\eta^{-1}\}$,
of Lorentz class (Proposition~17.4 and Theorem~17.11(ii),(iii) of \cite{ES17}); an invariant cross block vanishes when only one constituent is acted on, and
is a complex multiple of the invariant bilinear form under $\lG^\Delta$. This gives $2\cdot2=4$ and $2\cdot3=6$, the real and imaginary parts of the
products $z_k\cdot z_l$ of \cref{e5:prop:quad}, and $\so{6,\C}$ keeps only $q_{\rm tot}$, which gives $2$. A copy rotation acts on the copy index, and the
copy-rotation-invariant complex symmetric $2\times2$ matrices are the multiples of $\id$, which gives $2$ for the stabiliser. Under $\lG^{(2)}$ constituent~$1$ is
free ($12$ and $9$ by Theorem~17.11(i) of \cite{ES17}), and its cross blocks vanish because $\so3$ moves every vector of $\R^3$. Under the compact $\su2$
alone the identity of constituent~$2$ is also invariant, which adds one duality form. Under $\R C_v$ alone, $N=\id\otimes A+\varepsilon\otimes B$ with $A$
symmetric and $B$ antisymmetric, which gives $12+6$ and $9+9$. With no symmetry the classes are the forms anticommuting and commuting with $G_n$.
\end{proof}

The first, third and fourth rows restate Section~8.2 of \cite{RE} (``four under $\lG^{(1)}\oplus\lG^{(2)}$, six under $\lG^\Delta$ alone and two under
$\Cent(G_2)\cap\so{6,6}$''); the other rows, and the ladder, are fixed (exact ranks; e5-4a, e5-4b). Two observations follow. The seven-dimensional
stabiliser of the pair invariant fixes the inertia as much as the thirty-dimensional closure symmetry does: up to one complex scale. And, as at $n=3$, imposing
only what the binding root leaves of a gauge algebra enlarges the freedom to $14$ or $18$ real parameters in the Lorentz class, so the gauge parts do not
select the inertia.

\subsection{Two constituents and one rotor}
\label{e5:sec:rotor}

The closure on $\mathfrak s_2=\R\ltimes_{G_2}V_2$ with $H=h^2/(2I_0)+\sum_k\half p^{(k)\top}N_kp^{(k)}$ has $p^{(k)}=e^{-\theta G}p^{(k)}(0)$ for both
constituents, so $w_k=e^{2i\theta}w_k(0)$ and $I_0\ddot\theta=R_{\rm eff}\sin(2\theta+\varphi_{\rm eff})$. The rate is constant and nonzero iff
$R_{\rm eff}=0$ and $h(0)\neq0$, and then every frame moves on a circle of diameter $|N_kp^{(k)}(0)|/|\kappa|$. The balance set $\mathcal B=\{F=0\}$ has
codimension two away from $0$, its massive part is open and dense, and on it $|c_1|E_0(\pi^{(1)})=|c_2|E_0(\pi^{(2)})$ (Theorem~8.1 of \cite{RE}, restating
Proposition~17.18 of \cite{ES17}). With equal inertias the pairs on the null subspaces $V_\pm=\{p^{(2)}=\mp Gp^{(1)}\}$ are balanced (Proposition~8.4
of \cite{RE}). The following items are new.

\begin{proposition}[The rotor law with pair inertias]
\label{e5:prop:rotorlaw}
\begin{enumerate}[label=(\roman*),leftmargin=2.2em]
\item $\half p^{\!\top}Np=\Re\bigl((a+ib)w(p)\bigr)/(2m)$ for $N=(a\eta^{-1}+bG\eta^{-1})/m$, so block-diagonal Lorentz-class inertias give
$U(\theta)=\half\Re\bigl(e^{2i\theta}F(p(0))\bigr)$.
\item The Lorentz-class forms on $V_2^*$ invariant under $\lG^\Delta$ are exactly the real and imaginary parts of the three products $z_k\cdot z_l$ of the
momentum fields. For such an inertia $U=\half\Re\tr\bigl(C\,Q(\pi)\bigr)$ with $C$ complex symmetric, $Q(\pi(\tau))=e^{2i\theta}Q(\pi(0))$, and
\[
I_0\ddot\theta=R\sin(2\theta+\varphi),\qquad Re^{i\varphi}=\tr\bigl(C\,Q(\pi(0))\bigr).
\]
The balance condition is $\tr(CQ(\pi))=0$, and $C=\operatorname{diag}(c_1,c_2)$ recovers (i).
\end{enumerate}
\end{proposition}

\begin{proof}
(i) is Theorem~17.11(ii) of \cite{ES17} per constituent. (ii) The first statement is \cref{e5:thm:inertia} (row $\lG^\Delta$) with \cref{e5:prop:quad}(c). Every
$z(\pi^{(k)})$ turns as $e^{i\theta}$, so $Q\mapsto e^{2i\theta}Q$, $U(\theta)=\half R\cos(2\theta+\varphi)$, and $\dot h=-dU/d\theta$.
\end{proof}

Part (i) restates Theorem~17.11(ii) of \cite{ES17} (symbolic). Part (ii) restates the remark in Section~8.2 of \cite{RE}, which states the form and the
balance condition; the rotor law with cross terms is fixed here. Checks: exact span, and the law along a flow with a random $C$ to $10^{-9}$ (e5-5a).

\begin{proposition}[The balance surface at two constituents]
\label{e5:prop:Bgeom}
Let $c_1,c_2\neq0$.
\begin{enumerate}[label=(\roman*),leftmargin=2.2em]
\item In the data $(w_1,w_2)\in\C^2$ the balance condition $c_1w_1+c_2w_2=0$ is a complex line, of real codimension two. In the momentum data
$\R^{12}$, $\mathcal B$ has codimension two away from $0$. If exactly one coefficient vanishes, say $c_2=0$, then $\mathcal B=\{w_1=0\}\times\R^6$
still has codimension two, but forces constituent~$1$ to be null; if both vanish, $\mathcal B$ is everything.
\item $\mathcal B\setminus0$ is connected. For $c_1=c_2$ it is the nonzero null cone $\{Z\cdot Z=0,\ Z\neq0\}$ of $Z=(z(\pi^{(1)}),z(\pi^{(2)}))\in\C^6$,
a single orbit of $SO(6,\C)$, diffeomorphic to $\R_{>0}\times V_2(\R^6)$ with $V_2(\R^6)$ the Stiefel manifold of orthonormal two-frames. For $c_1\neq c_2$
it is the image of that cone under $(z_1,z_2)\mapsto(z_1,\rho^{-1}z_2)$, $\rho^2=c_2/c_1$.
\item The all-null part of $\mathcal B$ has dimension $8$. The massive part, on which both constituents are massive, has dimension $10$, and $w_1$ maps it
submersively onto $\C^*$, with fibre over $s$ the product of the two complex spheres $\{w(p^{(1)})=s\}$ and $\{w(p^{(2)})=-c_1s/c_2\}$.
\end{enumerate}
\end{proposition}

\begin{proof}
(i) is Theorem~8.1(iii) of \cite{RE}; the degenerate cases are direct. (ii) $c_1q_1+c_2q_2=c_1(z_1,\rho z_2)\cdot(z_1,\rho z_2)$. A nonzero null vector of
$\C^6$ is $u+iv$ with $u,v\in\R^6$, $|u|=|v|>0$ and $u\perp v$; $SO(6,\C)$ is transitive on these by Witt's theorem \cite{Witt37}, and $V_2(\R^6)$ is
connected. (iii) The null set $\{w=0\}\subset\R^6$ has dimension four (Theorem~8.1(iii) of \cite{RE}). On the massive part $w_1\neq0$ and $w_2=-c_1w_1/c_2\neq0$.
\end{proof}

Fixed, with (i) restated. Checks at $20$ random points each: $\operatorname{rank}dF=2$ for equal and unequal coefficients, orbit dimension $10$,
$\operatorname{rank}d(w_1,w_2)=4$ on the all-null part, and $\operatorname{rank}(dF,dw_1)=4$ on the massive part (e5-5f).

\begin{theorem}[Stability of the balance surface]
\label{e5:thm:stab}
Let $c_1,c_2\neq0$ and let $(\theta,h,p,x)$ solve the two-constituent closure with $\theta(0)=0$, $\kappa_0=h(0)/I_0$ and $\psi_0=\varphi_{\rm eff}$.
\begin{enumerate}[label=(\roman*),leftmargin=2.2em]
\item $R_{\rm eff}$ is a constant of motion for every datum. The balance surface and every level set $\{R_{\rm eff}=r\}$ are invariant, so the
surface is neither attracting nor repelling. A perturbation $\delta p$ of a balanced datum has imbalance $|dF(\delta p)|+O(|\delta p|^2)$, and since $dF$ is onto
every small imbalance occurs.
\item The rate is stable. With $\varepsilon=R_{\rm eff}/(I_0\kappa_0^2)$, $\kappa^2=\kappa_0^2[1+\varepsilon(\cos\psi_0-\cos\psi)]$ and
$|\kappa^2/\kappa_0^2-1|\le2\varepsilon$ for all $\tau$, and the rotor turns for all time when $\varepsilon<\half$.
\item The bounded frame orbits are not stable. In the rotating regime, during each period of the rate frame $k$ is displaced by
$\lambda\,e^{-\varphi_{\rm eff}G}N_kp^{(k)}(0)$, with $\lambda$ of Theorem~17.16 of \cite{ES17}, and its mean velocity is
\[
\bar v_k=\tfrac\varepsilon4\bigl(1-\varepsilon\cos\psi_0\bigr)\,e^{-\varphi_{\rm eff}G}N_kp^{(k)}(0)+O(\varepsilon^3),
\]
of first order in the imbalance. A perturbation off the surface therefore moves the frames by a distance that grows linearly in $\tau$, at a rate linear
in the perturbation. A perturbation along the surface, or of $h(0)$, keeps them on circles.
\item If $h(0)=0$, a balanced datum is an equilibrium of the rotor, and each frame moves uniformly, $x_k=x_k(0)+\tau N_kp^{(k)}(0)$, without bound.
Off the surface, unless $\psi_0\in\pi\Z$ (an equilibrium, at which the frames move uniformly), the rotor librates about $\psi=\pi$ with
$\Omega^2=2R_{\rm eff}/I_0$, amplitude $\xi_0\in(0,\pi)$ given by $\cos\xi_0=-\cos\psi_0$, and period $4\mathrm K(m)/\Omega$, $m=\sin^2(\xi_0/2)$,
which diverges as $R_{\rm eff}\to0$; $\mathrm K$ and $\mathrm E$ are the complete elliptic integrals in the parameter $m$ \cite{DLMF}. The
frames move with the mean velocity $\bigl(1-2\mathrm E(m)/\mathrm K(m)\bigr)e^{-\varphi_{\rm eff}G}N_kp^{(k)}(0)$. They drift, except at the one
amplitude with $2\mathrm E(m)=\mathrm K(m)$ ($m=0.8261$, $\xi_0=130.7^\circ$), at which every frame returns to its initial position after each
period.
\end{enumerate}
\end{theorem}

\begin{proof}
(i) $F(p(\tau))=e^{2i\theta}F(p(0))$ because every $w_k$ turns by $e^{2i\theta}$. (ii) Energy conservation, $H=\half I_0\kappa^2+\half R_{\rm eff}\cos\psi$, divided
by $\half I_0\kappa_0^2$; this is Corollary~17.15 of \cite{ES17} with $R_0/m$ replaced by $R_{\rm eff}$. (iii) Lorentz-class $N_k$ anticommute with $G$, so
$\dot x_k=e^{2\theta G}N_kp^{(k)}(0)$ with the common angle $\theta$, and the proof of Theorem~17.16 of \cite{ES17} applies to each frame. Expanding
$(1+\varepsilon(\cos\psi_0-\cos\psi))^{-1/2}$ to second order,
$\lambda=\frac{\pi\varepsilon}{4|\kappa_0|}\bigl(1-\tfrac32\varepsilon\cos\psi_0\bigr)+O(\varepsilon^3)$ and the period is
$\frac\pi{|\kappa_0|}\bigl(1-\half\varepsilon\cos\psi_0\bigr)+O(\varepsilon^2)$; their ratio gives $\bar v_k$. Along the surface $\varepsilon=0$, and a change
of $h(0)$ changes only the radius $|N_kp^{(k)}(0)|/(2|\kappa|)$ of the circles. (iv) At $h(0)=0$ and $R_{\rm eff}=0$, $\ddot\theta=0$ and $\theta\equiv0$, so
$\dot x_k=N_kp^{(k)}(0)$. Off the surface the rotor equation is the pendulum of Proposition~4.4 of \cite{RE} with $R_0/m$ replaced by $R_{\rm eff}$. With
$\psi=\pi+\xi$, $\dot x_k=-e^{-\varphi_{\rm eff}G}(\cos\xi+G\sin\xi)N_kp^{(k)}(0)$. The substitution $\sin(\xi/2)=\sin(\xi_0/2)\sin t$ of that
proposition gives $d\tau=dt/(\Omega\sqrt{1-m\sin^2t})$, so over one period $\int\sin\xi\,d\tau=(2\sqrt m/\Omega)\int_0^{2\pi}\sin t\,dt=0$ and
$\int\cos\xi\,d\tau=(4/\Omega)\int_0^{\pi/2}(1-2m\sin^2t)(1-m\sin^2t)^{-1/2}dt=(4/\Omega)\bigl(2\mathrm E(m)-\mathrm K(m)\bigr)$; divide by the
period.
\end{proof}

Item (i) restates Theorem~8.1(i) of \cite{RE}; (ii) restates Corollary~17.15 of \cite{ES17} for two constituents; (iii) and (iv) are fixed. The
second-order term in (iii) is new. Item (iv) sharpens the closing sentence of Proposition~17.18 of \cite{ES17}: on the balance surface the frames are
bounded only if $h(0)\neq0$, which is the qualifier ``constant and nonzero'' of Theorem~8.1(i) of \cite{RE}. Checks (DOP853, relative and absolute
tolerance $10^{-12}$; e5-5b to e5-5e): with equal inertias $a+ib=1+0.3i$, $m=I_0=1$, $h(0)=2$, and $p^{(2)}=-Gp^{(1)}$ drawn at random
($E_0(\pi^{(k)})=1.2156$ each), the rate stays constant to $10^{-11}$ over $\tau\in[0,60]$ and both frames lie on their circles (diameter $0.942$) to
$10^{-9}$. A generic balanced pair with unequal inertias ($m=1,2.5$; $a+ib=1+0.3i,\ 0.4-0.7i$) gives the same. Off the surface ($p^{(2)}$ perturbed by
$0.05$ along a fixed random direction, $R_{\rm eff}=0.391$, $\varepsilon=0.098$), $R_{\rm eff}$ and $w_1\overline{w_2}$ stay constant to $10^{-10}$, the
rotor obeys the pendulum law to $10^{-9}$, and the displacement per period is $\lambda e^{-\varphi_{\rm eff}G}N_kp^{(k)}(0)$ to relative $10^{-8}$. For
perturbations $10^{-3}$, $2\times10^{-3}$, $4\times10^{-3}$ the imbalance, the rate excursion and the mean frame velocity double with the perturbation (to
$2\%$). The mean velocity equals the quadrature value to $10^{-8}$, the first-order law to within $1.5\varepsilon$ and the second-order law to within
$\varepsilon^2$ (relative). With $h(0)=0$ the rotor librates with the period of (iv) to $10^{-8}$, and on the balance surface it stays at rest to
$10^{-12}$ while the frames move uniformly. Off the surface with $h(0)=0$ the displacement after one, two and three periods equals
the mean-velocity law of (iv) to relative $10^{-9}$ at four amplitudes, among them $\xi_0=130.7^\circ$, where it vanishes (e5-5i).

\begin{proposition}[Exchange of the constituents]
\label{e5:prop:exchange}
Let $\sigma(p^{(1)},p^{(2)})=(p^{(2)},p^{(1)})$.
\begin{enumerate}[label=(\roman*),leftmargin=2.2em]
\item $\sigma$ is an $\eta_2$-isometry commuting with $G_2$, so it preserves the Lie--Poisson bracket of $\mathfrak s_2$ and commutes with the rotor; it
preserves the equal-inertia Hamiltonian; as an element of $O(6,\C)$ it has determinant $-1$, so it lies outside the identity component of the closure
symmetry group. It exchanges $V_+$ and $V_-$.
\item With equal coefficients $R_{\rm eff}=2|c|\,|w_s|$ with $w_s=\half(w_1+w_2)$, the exchange-even part of $(w_1,w_2)$: a pair is balanced iff
$(w_1,w_2)$ is exchange-odd, $w_2=-w_1$.
\item The exchange eigenstates of the momenta, $p^{(2)}=\pm p^{(1)}$, both have $w_2=w_1$, so $F=(c_1+c_2)w_1$: unless $c_2=-c_1$ they are
balanced only when null. With equal inertias a massive balanced pair is never an exchange eigenstate of its momenta.
\end{enumerate}
\end{proposition}

\begin{proof}
(i) $\sigma$ permutes the copies; on $Z$ it is the permutation of three pairs of coordinates, of determinant $(-1)^3$. If $p^{(2)}=-Gp^{(1)}$, then
$\sigma p$ has second component $p^{(1)}=G(-Gp^{(1)})$. (ii) $R_{\rm eff}=|c|\,|w_1+w_2|$. (iii) $w$ is quadratic, so $F=c_1w_1+c_2w_1$.
\end{proof}

Fixed (e5-5g, e5-5j). Exchange parity of the duality invariants and exchange parity of the momenta are different things: the balanced pairs of $V_\pm$, a
magnetic and an electric constituent a quarter duality turn apart (Proposition~8.4(iii) of \cite{RE}), are exchanged with each other by $\sigma$ and are
not exchange eigenstates. With $c_1\neq c_2$, $\sigma$ is not a symmetry of the closure.

\begin{proposition}[What carries over from the rest closure]
\label{e5:prop:carry}
For one rotor on two constituents with Lorentz-class inertias:
\begin{enumerate}[label=(\roman*),leftmargin=2.2em]
\item (Theorem~4.1 of \cite{RE}.) A function of $p$ is a Casimir iff it is invariant under $p\mapsto e^{tG_2}p$. At $p\neq0$ the leaves are
two-dimensional, so there are $11$ local Casimirs, against $5$ at $n=1$, and the real and imaginary parts of $\bar z_az_b$, $a,b=1,\dots,6$, with $z\in\C^6$ the pair of
momentum fields, provide them. They include $E_6$, $E_0$ and $\mathbf S$ of each $\pi^{(k)}$ and the moduli $|Q_{kl}|$; no Hamiltonian moves rest energy
between the constituents (Corollary~8.3 of \cite{RE}).
\item (Propositions~4.3--4.8 of \cite{RE}.) The pendulum, its small oscillations, the energy budget and the multiscale criterion hold with $R_0/m$ replaced
by $R_{\rm eff}$, which depends on the data only through $F(p(0))$. What changes is Corollary~4.5(i) of \cite{RE}: a constant nonzero rate no longer
needs a null momentum field.
\item (Propositions~4.10 and~4.14 of \cite{RE}.) $\Gamma_{12}^2=\id$, and $\Gamma_{12}$ anticommutes with the twelve generators, so it implements
$e^{\pi G_2}=-\id$ on $S_2$. Since $(e^{\theta G_2})^{\!\top}\eta_2e^{\theta G_2}=\cos2\theta\,\eta_2+\sin2\theta\,\eta_2G_2$, the rotor lifts to
$\mathrm{Pin}(6,6)$ only at $\theta\in\pi\Z$.
\item (Proposition~4.9 of \cite{RE}.) The rotor turns both momentum fields at one rate, while the graded rest Hamiltonian of the tower has the single
frequency $(m_1^2+m_2^2)^{1/2}$ and the ungraded one the two frequencies $m_1$, $m_2$ (\cref{e1:prop:rest}(v)). A closure with two rates exists, with
the rotor acting by $\operatorname{diag}(\kappa_1,\kappa_2)\otimes G$, $\kappa_1\neq\kappa_2$ (Proposition~9.2 of \cite{RE}), and is admitted, not selected.
\item (Proposition~4.16 of \cite{RE}.) Under Assumption~4.15 of \cite{RE}, every leaf through the balance surface is flat, $U\equiv0$, so the
conditional ladder $\kappa_n=(n+\delta)\hbar/I_0$ holds exactly with both constituents massive, where at $n=1$ it needs null data.
\end{enumerate}
\end{proposition}

\begin{proof}
(i) The proof of Theorem~4.1 of \cite{RE} with $G_2$ in place of $G$; $zz^\dagger$ is a rank-one Hermitian matrix on $\C^6$, with $2\cdot6-1=11$ real
parameters. (ii) The equations for $(\theta,h)$ are those of one constituent with $R_0e^{i\varphi_0}/m$ replaced by $F(p(0))$. (iii) As in Proposition~4.14 of
\cite{RE}: a product of all twelve anticommuting generators anticommutes with each, and $G_2^{\!\top}\eta_2=\eta_2G_2$, $G_2^{\!\top}\eta_2G_2=-\eta_2$.
(iv) and (v) are direct.
\end{proof}

Items (i)--(iii) restate the cited results for two constituents (fixed there) and are fixed in their $n=2$ form (Jacobian rank $11$ and $5$, Lie--Poisson
rank $2$, exact Clifford identities; e5-5h). Item (iv) restates \cref{e1:prop:rest} and Proposition~9.2 of \cite{RE} and is constructed for the two-rate
closure. Item (v) is constructed and conditional on the assumption. The copy rotor of Remark~8.10 of \cite{RE}, a second rotor that exists because the
uniqueness argument of Theorem~17.7 of \cite{ES17} fails under $\lG^\Delta$, and the two-rate rotor of (iv) are the structures of the rest closure that exist only from $n=2$ on.

\subsection{The ladder: one, two and three constituents}
\label{e5:sec:ladder}

\begin{table}[htbp]
\centering
\footnotesize
\begin{tabular}{@{}>{\raggedright\arraybackslash}p{3.4cm}>{\raggedright\arraybackslash}p{2.9cm}>{\raggedright\arraybackslash}p{4.3cm}>{\raggedright\arraybackslash}p{3.9cm}@{}}
\toprule
 & $n=1$ & $n=2$ & $n=3$\\
\midrule
closure symmetry $\Cent(G_n)$ & $\lG\cong\so{3,\C}$, $6$ & $\so{6,\C}$, $30$ & $\so{9,\C}$, $72$\\
$\Cent(G_n)\cap\lint^{(1)}$ & $0$ & $0$ & $0$\\
commuting with $G_n$ and $\lint^{(1)}$ & $0$ & $\lG^{(2)}$, $6$ & $\so{6,\C}$ of $2,3$, $30$\\
$S_n\otimes\C$ under $\lG^{(1)}$ & $2\,(\half,\half)$ & $16\,(\half,\half)$ & $128\,(\half,\half)$\\
Pati--Salam algebra & no & no & in $\so8\oplus\so6$\\
binding-compatible compact part of the spectator algebra & $0$ & $\su2$, $3$; vector-like, isospin $0,1$ & $\so6$ of $2,3$, $15$; contains $\so4_\Delta$\\
the same for the diagonal connection algebra & $0$ & $\R C_v$, $1$; no fermions & copy rotations $\so3$, $3$\\
under the projectors & --- & --- & $\Pfam$: $\uu1^2$; $P_*$: $\so4_\Delta\oplus\R\,\mathrm{rot}(8,9)$, chiral\\
anomalies of what survives & --- & none & $\Pfam$: $\mathcal A(H^3)=-12\alpha\beta^2$; $P_*$: free on $\so4_\Delta$\\
inertia, closure symmetry & $(2,0)$ & $(2,0)$ & $(2,0)$\\
inertia, per-constituent $\lG$ & $(2,0)$ & $(4,0)$ & $(6,0)$\\
inertia, binding-compatible part & --- & $(14,10)$ spectator; $(18,18)$ diagonal & $(18,13)$ $\Pfam$; $(8,6)$ $P_*$\\
inertia, none & $(12,9)$ & $(42,36)$ & $(90,81)$\\
balance surface & $\{w=0\}$: null data only, dim $4$ & codim $2$, connected; massive part dense, dim $10$ & codim $2$, dim $16$\\
\bottomrule
\end{tabular}
\caption{The closure along the ladder. Dimensions are real; the inertia entries are (Lorentz class, duality class). The $n=3$ column restates
Theorems~17.19--17.24 and Section~10.7 of \cite{ES17} and Theorem~8.1(iii) of \cite{RE}, except the spectator and diagonal rows, which
are computed here (e5-6a).}
\label{e5:tab:ladder}
\end{table}

\begin{table}[htbp]
\centering
\footnotesize
\begin{tabular}{@{}>{\raggedright\arraybackslash}p{8.0cm}>{\raggedright\arraybackslash}p{4.7cm}>{\raggedright\arraybackslash}p{2.1cm}@{}}
\toprule
statement & class & checks\\
\midrule
\cref{e5:prop:so6C}: $\Cent(G_2)=\so{6,\C}$, paired rotations and boosts, $R-iS$, compact part $\so6$ & restated (RE Prop.~7.8(i); ES17 Thm~17.19 at $n=2$) & e5-1a\\
\cref{e5:prop:quad}: quadratic invariants $q_{\rm tot}$; $q_1,q_2$; $q_1,q_2,z_1\cdot z_2$ & inherited (invariant theory); dimensions fixed & e5-1b\\
\cref{e5:thm:pair}(i): $M_{\rm pair}$ through the Gram entries; diagonal against per-constituent invariance & fixed (partly restated, RE Thm~8.6) & e5-1e\\
\cref{e5:thm:pair}(ii): stabiliser of $M_{\rm pair}$ is $\lG^\Delta\oplus\R C_v$, dimension $7$ & fixed (numerical rank) & e5-1c\\
\cref{e5:thm:pair}(iii),(iv): $\so{6,\C}$ does not preserve $M_{\rm pair}$; range $[\half|c|,\infty)$ on each orbit & fixed (Ex.~8.8 of RE restated; transitivity inherited) & e5-1d\\
\cref{e5:prop:nocommon}: no common Lorentz algebra; $\lG^{(2)}$; $16(\half,\half)$; integer spin under $\lG^\Delta$ & fixed (restated: RE Props.~7.8(iv), 9.1(iii); ES17 Thm~17.20) & e5-2a to e5-2c\\
\cref{e5:prop:noPS}: no Pati--Salam algebra at $n=2$ & obstructed (dimension count) & e5-3a, 3e, 3h, 3i\\
\cref{e5:prop:spectator}: binding-compatible $\su2$; kinetic condition; integer isospin, vector-like, anomaly free & fixed (Weyl block restated, ES17 Prop.~10.12; branching inherited) & e5-3a to e5-3d, 3j\\
\cref{e5:prop:diagonal}: $\R C_v$; no canonical kinetic kernel; content, vector-like & fixed ((i) restated, RE Props.~9.1(ii), 7.3) & e5-3e to e5-3g\\
\cref{e5:prop:mixed}: mixed rung, nothing survives $G_2$; covariant rotors $G^{(2)},J_1$ & fixed (restated: ES17 Props.~13.16, 13.17); rotor $G^{(2)}$ constructed & e5-3h\\
binding actions: $\lG^{(2)}$ and $\so{2,\C}$ survive & restated (RE Prop.~7.8(iii), Table~8) & e5-3i\\
\cref{e5:thm:inertia}: inertia freedom at $n=2$; $(2,0)$ at $n=1,2,3$ & fixed (rows 1, 3, 4 restated, RE Sec.~8.2) & e5-4a, e5-4b\\
\cref{e5:prop:rotorlaw}: rotor law with pair inertias & fixed ((i) restated, ES17 Thm~17.11(ii)) & e5-5a\\
\cref{e5:prop:Bgeom}: codimension two, connected, one orbit, parts & fixed ((i) restated, RE Thm~8.1(iii)) & e5-5f\\
\cref{e5:thm:stab}: neutral surface, stable rate, secular frames, $h(0)\neq0$ needed, libration drift & fixed ((i), (ii) restated) & e5-5b to e5-5e, 5i\\
\cref{e5:prop:exchange}: exchange, $\det\sigma=-1$, balanced iff $w$ exchange-odd & fixed & e5-5g, 5j\\
\cref{e5:prop:carry}: Casimirs, pendulum, half turn, rates, ladder at $n=2$ & (i)--(iii) restated (RE Sec.~4), fixed at $n=2$; (iv), (v) constructed & e5-5h\\
\cref{e5:tab:ladder}: the ladder $n=1,2,3$ & fixed; $n=3$ mostly restated (ES17 Sec.~17.6) & e5-6a, 6b, 2a, 4b\\
\cref{e5:sec:ladder}: answer & interpretive (drawn from fixed statements) & ---\\
\bottomrule
\end{tabular}
\caption{Claims of \cref{e5:sec:main}, their classes and their checks, which are \texttt{PASS} lines of the script
\texttt{verify/e5\_closure\_two.py}.}
\label{e5:tab:claims}
\end{table}

The new entries of \cref{e5:tab:ladder} at $n=1$ and $n=3$ are fixed (e5-6a). At $n=1$, $\Cent(\lint)\cap\so{3,3}=\R J$ meets $\Cent(G)$ in $0$, so the
binding root leaves no gauge algebra on one constituent. At $n=3$ the compact spectator algebra meets $\Cent(G_3)$ in the $15$ paired rotations of
constituents $2$ and $3$, which contain the $\so4_\Delta$ of Theorem~17.21 of \cite{ES17}, and the diagonal connection algebra meets it in the three copy
rotations. At every rung the binding-compatible parts are $\so{3n-3}$, the paired rotations of constituents $2,\dots,n$, and $\so n$, the copy
rotations: the proofs of \cref{e5:prop:spectator}(i) and \cref{e5:prop:diagonal}(i) apply verbatim (fixed; exact at $n=4$, e5-6b). The balance condition $F=0$ has codimension two at every rung, because $dF$
has rank two at every nonzero point. At $n=1$ it is the null cone and carries no massive datum.

\paragraph{Answer.} \emph{Is $n=2$ a closure rung in its own right?} Yes, in a precise sense (interpretive, drawn from the fixed statements cited here):
every structure of the three-constituent analysis of \cite{ES17} except the two projectors, which exist only at $n=3$, has a computed $n=2$
analogue, and four of them first appear at $n=2$. The balance surface carries massive constituents at
an exactly constant rate (Theorem~8.1 of \cite{RE}). The closure symmetry mixes the constituents and does not see their pair invariant
(\cref{e5:thm:pair}). The closure symmetry compatible with the matter Lorentz algebra is nonzero, and is itself a Lorentz algebra (\cref{e5:prop:nocommon}).
And the binding root leaves a compact gauge algebra, non-abelian under the spectator action (\cref{e5:prop:spectator}). \emph{What $n=2$ adds to $n=1$} is
therefore a constant rate with mass, a gauge part that survives the binding root, a pair invariant and an exchange symmetry. \emph{What it lacks}, and
\emph{what $n=3$ adds}, is everything chiral. At $n=2$ every binding-compatible gauge part acts vector-like, with zero anomaly coefficients and integer
isospin, and no Pati--Salam algebra fits. At $n=3$ the projectors $\Pfam$ and $P_*$ exist, the surviving $\so4_\Delta$ acts on inequivalent halves
with doublets, the abelian charges of $\Pfam$ carry nonzero anomaly coefficients, and fermions exist under the diagonal action. The inertia question does
not distinguish the rungs: the closure symmetry fixes the inertia up to one complex scale at $n=1,2,3$ alike, and the gauge parts never do.

For the question of this note the closure gives a definite but modest answer (interpretive, drawn from the fixed statements above). The closure at $n=2$
distinguishes two constituents only through structures it does not select. Equal inertias and the closure symmetry see one field strength in $\C^6$. A
pair invariant needs a diagonal Lorentz action and the copy rotation. Per-constituent rest energies need the copies to be distinguished, and then nothing in
the closure moves rest energy between them (Proposition~8.9 and Corollary~8.3 of \cite{RE}). What two constituents can do that one cannot is to run at
an exactly constant rate with both momentum fields massive. That holds on a set of codimension two which is invariant but whose frames are secularly
unstable, which is closer to a balanced pair carried by the dynamics than to two particles produced by it.
